\documentclass[10pt]{iopart}

\usepackage{epsfig,graphicx}
\usepackage{algpseudocode}
\usepackage{amsfonts}
%\usepackage{amsmath}
\usepackage{amsbsy}
\usepackage{dsfont}
\usepackage{color}
\usepackage{comment}
\usepackage{float}

\usepackage{hyperref}
\hypersetup{
    colorlinks,
    citecolor=black,
    filecolor=black,
    linkcolor=black,
    urlcolor=black
}

\usepackage{amsthm}

% for using amsmath
%\expandafter\let\csname equation*\endcsname\relax
%\expandafter\let\csname endequation*\endcsname\relax
%\usepackage{amsmath}

\newcommand{\new}{\textcolor{blue}}
\newcommand{\Cyan}{\textcolor{cyan}}
\newcommand{\old}{\textcolor{red}}


\newcommand{\sgn}{\textrm{sgn}}

\newcommand{\1}{\mathbf{1}}
\newcommand{\nullv}{\mathbf{0}}
\newcommand{\Ind}{\mathds{1}}

\newcommand{\Prob}{\mathrm{P}}

\newcommand{\xiv}{  \ensuremath{\boldsymbol{\xi}}   }

\newcommand{\W}{\mathrm{W}}

\newcommand{\Z}{\mathbb{Z}}
\newcommand{\N}{\mathbb{N}}
\newcommand{\mv}{\ensuremath{\mathbf{m}}}
\newcommand{\cv}{\ensuremath{\mathbf{c}}}
\newcommand{\sv}{\ensuremath{\mathbf{s}}}
\newcommand{\Sm}{\ensuremath{\mathbf{S}}}
%\newcommand{\etav}{\ensuremath{\boldsymbol{\eta}}}
\newcommand{\etav}{  \ensuremath{\boldsymbol{\eta}}   }
\newcommand\at[2]{\left.#1\right|_{#2}}
\newcommand{\E}{\mathbb{E}}
\newcommand{\Var}{\mathrm{Var}}

\newtheorem{prop}{Proposition}[section]
\newtheorem{theorem}{Theorem}[section]
\newtheorem{cor}{Corollary}[section]
\newtheorem{definition}{Definition}[section]
\newtheorem{assumption}{Assumption}[section]


\begin{document}

\title[Stochastic recurrence equation]{Analysis of stochastic recurrence equation}
%


\author{Alexander Mozeika}
\address{$^1$Status.im}


\ead{alexander.mozeika@status.im}




\date{\today}

%\maketitle

\begin{abstract}
We consider...  
\end{abstract}

\tableofcontents

\newpage
%%%%%%%%%%%%%%%%%%%%%%%%%%%%%%%%%%%%%%%%%%%%%%%%%%%%%%%%%%%%%%%%%%%%%%%%%%%%%%%


\section{Statistical properties of the leader election process \label{section:election-process}}
Start with the leader lottery equation:

\begin{equation}
\phi(\alpha)=1-(1-f)^\alpha,
\end{equation}
Where $\alpha\in[0,1]$ is relative stake and $f\in[0,1]$ is the \emph{active slots coefficient}.
To estimate the number of leaders per slot, we can look at the system evolved over time:

\begin{comment}
\begin{matrix}
    & \tau_0 & \tau_1 & \tau_2 & \tau_3 & \cdots \\
n_0 & 0 & 0 & 1 & 0 & \cdots \\
n_1 & 1 & 0 & \cdots \\
n_2 & 0 & 0 & \cdots \\
n_3 & 0 & 1 & \cdots \\
\vdots & \vdots & \vdots \\
n_N & 1 & 0 & \cdots\\
    & \vec{s^0} & \vec{s^1} & \vec{s^2} & \cdots
\end{matrix}
\end{comment}

\begin{center}
\begin{tabular}{ c c c c c}
         & $\tau_0$     & $\tau_1$      & $\tau_2$      & $\cdots$ \\ 
   $n_0$ & 0            & 0             & 1             & $\cdots$ \\
   $n_1$ & 1            & 0             & 0             & $\cdots$ \\
   $n_2$ & 0            & 0             & 0             & $\cdots$ \\
$\vdots$ & $\vdots$     & $\vdots$      & $\vdots$      &          \\
   $n_N$ & 1            & 0             & 0             & $\cdots$ \\
         & $\sv(0)$  & $\sv(1)$   & $\sv(2)$   & $\cdots$ \\
\end{tabular}
\end{center}

A $1$ in slot $\tau_i$ for node $n_j$ means $n_j$ was a leader for slot $\tau_i$, conversely, a $0$ means the corresponding node was not a leader in that slot.

The vector $\sv(t)$ gives a snapshot of all leaders/non-leaders for slot $\tau_t$.

With the system model in place, we can derive the distribution for $\sv(t)$:

\begin{eqnarray}
\Prob(\sv(t))&=&\prod_{i=1}^N \Prob(s_i(t))\nonumber\\
&=&\prod_{i=1}^N\left[\phi(\alpha_i)\,\delta_{1;s_i(t)}+(1-\phi(\alpha_i))\,\delta_{0;s_i(t)}\right]\label{def:Prob-s},
\end{eqnarray}
%
Where $\alpha_i\geq0$ is the relative stake of node $i$ and $\sum_{i=1}^N\alpha_i=1$. We note that $\alpha_i=w_i/\sum_{j=1}^Nw_j$, where $w_i>0$
 is stake of node $i$ and $\sum_{j=1}^Nw_j$ is the \emph{total} stake.  With this we can compute the average of random variable $\sum_{i=1}^N s_i(t)$, i.e. the \new{average} number of leaders in slot $t$, as follows 
%
\begin{eqnarray}
%
\left\langle\sum_{i=1}^N s_i(t)\right\rangle_{\sv(t)}&=& \sum_{\sv(t)}\Prob(\sv(t))\sum_{i=1}^N s_i(t) \nonumber\\
&=&\sum_{i=1}^N \sum_{\sv(t)} \Prob(\sv(t))s_i(t)  \nonumber\\
&=&\sum_{i=1}^N \sum_{s_i(t)} \Prob(s_i(t))s_i(t)\nonumber \\
&=&\sum_{i=1}^N\phi(\alpha_i)
%
\end{eqnarray}
%
Which intuitively makes sense, the expected number of leaders is the sum of running the lottery function for every node $n_i$. Thus the average is given by 
%
\begin{eqnarray}
%
\left\langle\sum_{i=1}^N s_i(t)\right\rangle_{\sv(t)}
&=&\sum_{i=1}^N\phi(\alpha_i)\label{eqn:aver}.
%
\end{eqnarray}
%

We also would like to know the variance 
%
\begin{eqnarray}
%
\mathrm{Var}\left[\sum_{i=1}^N s_i(t)\right]&=&\left\langle\left(\sum_{i=1}^N s_i(t)\right)^2\right\rangle_{\sv(t)} -\left\langle\sum_{i=1}^N s_i(t)\right\rangle_{\sv(t)}^2 \label{def:Var}
%
\end{eqnarray}
%
which  can be computed as follows. First, we consider the second moment  
%
\begin{eqnarray}
%
\left\langle\left(\sum_{i=1}^N s_i(t)\right)^2\right\rangle_{\sv(t)}&=&\sum_{i=1}^N\sum_{j=1}^N \sum_{\sv(t)} \Prob(\sv(t))s_i(t)s_j(t)  \nonumber\\
&=&\sum_{i=1}^N\sum_{j=1}^N \left\langle s_i(t)s_j(t)\right\rangle_{\sv(t)}
%
\end{eqnarray}
%
Now the average 
%
\begin{eqnarray}
%
\left\langle s_i(t)s_j(t)\right\rangle_{\sv(t)}&=&\delta_{i;j} \sum_{s_i(t)} \Prob(s_i(t))s_i(t)\nonumber \\
&& +(1-\delta_{i;j})\sum_{s_i(t)} \Prob(s_i(t))s_i(t)\sum_{s_j(t)} \Prob(s_j(t))s_j(t)\nonumber\\
&=&\delta_{i;j} \phi(\alpha_i) +(1-\delta_{i;j})\phi(\alpha_i)\phi(\alpha_j)
%
\end{eqnarray}
%
and hence 
%
\begin{eqnarray}
%
\left\langle\left(\sum_{i=1}^N s_i(t)\right)^2\right\rangle_{\sv(t)}&=&\sum_{i=1}^N\sum_{j=1}^N \Big[\delta_{i;j} \phi(\alpha_i)\nonumber\\
&&+(1-\delta_{i;j})\phi(\alpha_i)\phi(\alpha_j) \Big] \nonumber\\
&=&\sum_{i=1}^N\phi(\alpha_i)+\left(\sum_{i=1}^N \phi(\alpha_i)\right)^2\nonumber\\
&& -\sum_{i=1}^N\phi^2_f(\alpha_i).
%
\end{eqnarray}
%
Second, using above result  and the average  (\ref{eqn:aver}) in the definition  (\ref{def:Var}) gives us 
%
\begin{eqnarray}
%
\mathrm{Var}\left[\sum_{i=1}^N s_i(t)\right]
&=&\sum_{i=1}^N\phi(\alpha_i)+\left(\sum_{i=1}^N \phi(\alpha_i)\right)^2 -\sum_{i=1}^N\phi^2_f(\alpha_i)\nonumber\\
%
&&-\left(\sum_{i=1}^N \phi(\alpha_i)\right)^2\nonumber\\
%
&=&\sum_{i=1}^N\left[\phi(\alpha_i)-\phi^2_f(\alpha_i)\right]
%
\end{eqnarray}
%
and hence 
%
\begin{eqnarray}
%
\mathrm{Var}\left[\sum_{i=1}^N s_i(t)\right]&=&\sum_{i=1}^N\phi(\alpha_i)\left[1-\phi(\alpha_i)\right]\label{eqn:Var}
%
\end{eqnarray}
%



\section{Stochastic recurrence equation}

For $h>0$ we consider the equation 
%
\begin{eqnarray}
%
    D_{\ell+1}[\Sm(\ell)]&=& D_{\ell}[\Sm(\ell-1)]-h\, D_{\ell}[\Sm(\ell-1)]\nonumber\\
    &&\times\left(f-\frac{1}{T}\sum_{t=1}^T\1\left[\sum_{i=1}^N s^{\ell}_i(t)\ge1\right]\right)\label{eq:SRE},
    %
\end{eqnarray}
%\Sm(\ell)=\{\sv^\ell(t)\}
where in above $\Sm(\ell)\equiv \{\sv^{\ell}(1),\ldots,\sv^{\ell}(T)\}$ with $\sv^{\ell}(t)\in\{0,1\}^N$. The distribution of $\Sm(\ell)$ is given by 
%
\begin{eqnarray}
\Prob\left(\Sm(\ell)\vert D_{\ell}[\Sm(\ell-1)]\right)&=&\prod_{t=1}^T\Prob(\sv^\ell(t)\vert D_{\ell}[\Sm(\ell-1)])\nonumber\\
%
&=&\prod_{t=1}^T\prod_{i=1}^N \Prob(s^\ell_i(t)\vert D_{\ell}[\Sm(\ell-1)])\label{def:Prob-S},
\end{eqnarray}
%
where 
\begin{eqnarray}
\Prob(s_i(t)\vert D)&=&\phi(w_i/D)\,\delta_{1;s_i(t)}+(1-\phi(w_i/D))\,\delta_{0;s_i(t)}.
\end{eqnarray}
%
We note $\Prob(s_i(t)\vert D)$ can be written in the form
%
\begin{eqnarray}
\Prob(s_i(t)\vert D)&=&\frac{\rme^{s_i(t)\,h(w_i/D)}}{1+\rme^{h(w_i/D)}},
\end{eqnarray}
%
where  
%
\begin{eqnarray}
h(w_i/D)&=&\log\left(\frac{\phi(w_i/D)}{1-\phi(w_i/D)}\right).
\end{eqnarray}
%
We note that $1+\rme^{h(w_i/D)}=\frac{1}{1-\phi(w_i/D)}$. 

For $D_0$  independent  from $\Sm(0)$ above gives us 
%
\begin{eqnarray}
%
    D_{1}[\Sm(0)]&=&D_{0}-h\,D_{0} \left(f-\frac{1}{T}\sum_{t=1}^T\1\left[\sum_{i=1}^N s^0_i(t)\ge1\right]\right).
    %
\end{eqnarray}
%

For $\ell=1$ we obtain 
%
\begin{eqnarray}
%
    D_{2}[\Sm(1)]&=& D_{1}[\Sm(0)]-h \,D_{1}[\Sm(0)]\left(\!f-\frac{1}{T}\sum_{t=1}^T\1\left[\sum_{i=1}^N s^1_i(t)\ge1\right]\right),
    %
\end{eqnarray}
%
where
%
\begin{eqnarray}
\Prob\left(\Sm(1)\vert D_{1}[\Sm(0)]\right)&=&\prod_{t=1}^T\Prob(\sv^1(t)\vert D_{1}[\Sm(0)]).
\end{eqnarray}
%
The joint probability of $\Sm(1)$ and $\Sm(0)$ is given by 
%
\begin{eqnarray}
\Prob\left[\Sm(1),\Sm(0)\right]&=&\Prob\left(\Sm(1)\vert D_{1}[\Sm(0)]\right)\Prob(\Sm(0)),
\end{eqnarray}
%
where 
%
\begin{eqnarray}
\Prob_0(\Sm(0))&=&\Prob\left(\Sm(0)\vert D_0\right).
\end{eqnarray}
%
Furthermore, for $L\ge1$ we have the joint probability
%
\begin{eqnarray}
\Prob\left[\Sm(L),\ldots,\Sm(0)\right]&=&\Prob_0(\Sm(0))\prod_{\ell=1}^L\Prob\left[\Sm(\ell)\vert D_{\ell}[\Sm(\ell-1)]\right]\label{def:prob-path}.
\end{eqnarray}
%
Above is the probability of the path $\Sm(0)\to\cdots\to\Sm(L)$. Now we can define the joint probability 
%
\begin{eqnarray}
\Prob\left[\Sm(L),\ldots,\Sm(0);D_{L},\ldots,D_{1}\vert D_{0}\right]&=&\Prob_0(\Sm(0))\prod_{\ell=1}^L\Prob\left[\Sm(\ell)\vert D_{\ell}\right]\nonumber\\
&&\times\delta\left(D_{\ell}-D_{\ell}[\Sm(\ell-1)]\right)\label{def:prob-path-joint}
\end{eqnarray}
%
from which follows the probability distribution 
%
\begin{eqnarray}
\Prob\left[D_{L},\ldots,D_{1}\vert D_{0}\right]&=&\sum_{\Sm(0)}\cdots\sum_{\Sm(L)}\Prob_0(\Sm(0))\prod_{\ell=1}^L\Prob\left[\Sm(\ell)\vert D_{\ell}\right]\nonumber\\
&&~~~~~~~~~\times\delta\left(D_{\ell}-D_{\ell}[\Sm(\ell-1)]\right)\label{def:prob-path-D}.
\end{eqnarray}
%
For the above probability distribution  we can show the following result. 

\subsection{The binomial reduction}

\begin{prop}
%
\label{prop:D-distr}
The stochastic recurrence equation (\ref{eq:SRE}) is equivalent to the equation 
%
\begin{eqnarray}
%
D_{\ell+1}&=& D_{\ell}- h\,D_{\ell}\left(f- 1+\frac{k_{\ell}}{T}\right),
%
\end{eqnarray}
%
where $k_{\ell}\in\{0,1,\ldots,T\}$ is a random variable from the binomial distribution
%
\begin{eqnarray}
\Prob[k_{\ell}\vert D_{\ell}]={T\choose k}\left[1-\phi\left(\frac{\sum_{i=1}^Nw_i}{D_{\ell}}\right)\right]^k\phi^{T-k}\left(\frac{\sum_{i=1}^Nw_i}{D_{\ell}}\right).
\end{eqnarray}
% 

\end{prop}
\begin{proof}
Using the Fourier representation 
%
\begin{eqnarray}
\delta\left(D_{\ell}-D_{\ell}[\Sm(\ell-1)]\right)=\int_{-\infty}^{\infty}\frac{\rmd\hat{D}_\ell}{2\pi/T}\rme^{\rmi T\hat{D}_\ell\left(D_{\ell}-D_{\ell}[\Sm(\ell-1)]\right)}.
\end{eqnarray}
%
in (\ref{def:prob-path-D}) we obtain
%
\begin{eqnarray}
\Prob\left[D_{L},\ldots,D_{1}\vert D_{0}\right]&=&\sum_{\Sm(0)}\cdots\sum_{\Sm(L)}\Prob_0(\Sm(0))\prod_{\ell=1}^L\Prob\left[\Sm(\ell)\vert D_{\ell}\right]\nonumber\\
&&~~~~~~~~~\times\int_{-\infty}^{\infty}\frac{\rmd\hat{D}_\ell}{2\pi/T}\rme^{\rmi T\hat{D}_\ell\left(D_{\ell}-D_{\ell}[\Sm(\ell-1)]\right)}\nonumber\\%\label{def:prob-path-D},
%
&=&\left\{\prod_{\ell=1}^L\int_{-\infty}^{\infty}\frac{\rmd\hat{D}_\ell}{2\pi/T}\rme^{\rmi T\hat{D}_\ell D_{\ell}}\right\}\nonumber\\
&&\times\sum_{\Sm(0)}\cdots\sum_{\Sm(L)}\Prob_0(\Sm(0))\prod_{\ell=1}^L\Prob\left[\Sm(\ell)\vert D_{\ell}\right]\nonumber\\
&&~~~~~~~~~\times\rme^{-\rmi T\hat{D}_\ell D_{\ell}[\Sm(\ell-1)]}
%
\end{eqnarray}
%
Let us consider the sum 
%
\begin{eqnarray}
&&\sum_{\Sm(0)}\cdots\sum_{\Sm(L)}\Prob_0(\Sm(0))\prod_{\ell=1}^L\Prob\left[\Sm(\ell)\vert D_{\ell}\right]\rme^{-\rmi T\hat{D}_\ell D_{\ell}[\Sm(\ell-1)]}\nonumber\\
%
&&=\sum_{\Sm(0)}\cdots\sum_{\Sm(L)}\Prob_0(\Sm(0))\prod_{\ell=1}^L\Prob\left[\Sm(\ell)\vert D_{\ell}\right]\nonumber\\
%
&&\times\prod_{\ell=0}^{L-1}\rme^{-\rmi T\hat{D}_{\ell+1} D_{\ell+1}[\Sm(\ell)]}
%
\end{eqnarray}
%
First we consider
%
\begin{eqnarray}
&&\sum_{\Sm(0)}\Prob_0(\Sm(0))\,\rme^{-\rmi T\hat{D}_{1} D_{1}[\Sm(0)]}\nonumber\\
&&~=\rme^{-\rmi T\hat{D}_{1} \left\{D_{0}-h\,D_{0}f \right\}}\nonumber\\
%
&&~~~\times\sum_{\Sm(0)}\Prob_0(\Sm(0))\,\rme^{-\rmi h\,\hat{D}_{1} \,D_{0} \sum_{t=1}^T\1\left[\sum_{i=1}^N s_i^0(t)\ge1\right]}\nonumber\\
%
&&~=\rme^{-\rmi T\hat{D}_{1} \left\{D_{0}-h\,D_{0}f \right\}}\nonumber\\
%
&&~~~\times\sum_{\Sm(0)}\prod_{t=1}^T\left\{\prod_{i=1}^N\frac{\rme^{s^0_i(t)\,h(w_i/D_0)}}{1+\rme^{h(w_i/D_0)}}\right\}\,\rme^{-\rmi h\,\hat{D}_{1} \,D_{0} \1\left[\sum_{i=1}^N s_i^0(t)\ge1\right]}\nonumber\\
%
&&~=\rme^{-\rmi T\hat{D}_{1} \left\{D_{0}-h\,D_{0}f \right\}}\nonumber\\
%
&&~~~\times\prod_{t=1}^T\sum_{\sv^0(t)}\left\{\prod_{i=1}^N\frac{\rme^{s^0_i(t)\,h(w_i/D_0)}}{1+\rme^{h(w_i/D_0)}}\right\}\,\rme^{-\rmi h\,\hat{D}_{1} \,D_{0} \1\left[\sum_{i=1}^N s_i^0(t)\ge1\right]}\nonumber\\
%
&&~=\rme^{-\rmi T\hat{D}_{1} \left\{D_{0}-h\,D_{0}f \right\}}\nonumber\\
%
&&~~~\times\prod_{t=1}^T\Bigg[\left\{\prod_{i=1}^N\frac{1}{1+\rme^{h(w_i/D_0)}}\right\}\nonumber\\
%
&&~~~+\sum_{\sv^0(t)}\1\left[\sum_{i=1}^N s_i^0(t)\ge1\right]\left\{\prod_{i=1}^N\frac{\rme^{s^0_i(t)\,h(w_i/D_0)}}{1+\rme^{h(w_i/D_0)}}\right\}\,\rme^{-\rmi h\,\hat{D}_{1} \,D_{0} }\Bigg]\nonumber\\
%
&&~=\rme^{-\rmi T\hat{D}_{1} \left\{D_{0}-h\,D_{0}f \right\}}\nonumber\\
%
&&~~~\times\prod_{t=1}^T\left[\left\{\prod_{i=1}^N\frac{1}{1+\rme^{h(w_i/D_0)}}\right\}+\left\{1-\prod_{i=1}^N\frac{1}{1+\rme^{h(w_i/D_0)}}\right\}\,\rme^{-\rmi h\,\hat{D}_{1} \,D_{0} }\right]\nonumber\\
%
&&~=\rme^{-\rmi T\hat{D}_{1} \left\{D_{0}-h\,D_{0}f \right\}}\nonumber\\
%
&&~~~\times\left[\prod_{i=1}^N\left[1-\phi(w_i/D_0)\right]+\left\{1-\prod_{i=1}^N\left[1-\phi(w_i/D_0)\right]\right\}\,\rme^{-\rmi h\,\hat{D}_{1} \,D_{0} }\right]^T\nonumber\\
%
&&~=\rme^{-\rmi T\hat{D}_{1} \left\{D_{0}-h\,D_{0}f \right\}}\nonumber\\
%
&&~~~\times\left[\Prob(\mathbf{0}\vert D_0)+\left(1-\Prob(\mathbf{0}\vert D_0)\right)\,\rme^{-\rmi h\,\hat{D}_{1} \,D_{0} }\right]^T\nonumber\\
%
&&~=\rme^{-\rmi T\hat{D}_{1} \left\{D_{0}-h\,D_{0}f \right\}}\sum_{k=0}^T{T\choose k}\Prob^k(\mathbf{0}\vert D_0)\left[1-\Prob(\mathbf{0}\vert D_0)\right]^{T-k}\,\rme^{-\rmi h\,\hat{D}_{1} \,D_{0} (T-k)},
%
\end{eqnarray}
%
where in above we defined the prob.
%
\begin{eqnarray}
\Prob(\mathbf{0}\vert D_0)=\prod_{i=1}^N\left[1-\phi(w_i/D_0)\right]. 
\end{eqnarray}
%
We note that 
\begin{eqnarray}
1-\Prob(\mathbf{0}\vert D_0)=\phi\left(\sum_{i=1}^Nw_i/D_0\right). 
\end{eqnarray}
%


Hence from above follows that 
%
\begin{eqnarray}
&&\sum_{\Sm(0)}\Prob\left(\Sm(0)\vert D_0\right)\,\rme^{-\rmi T\hat{D}_{1} D_{1}[\Sm(0)]}\nonumber\\
&&~=\rme^{-\rmi T\hat{D}_{1} \left\{D_{0}-h\,D_{0}f \right\}}\sum_{k=0}^T{T\choose k}\Prob^k(\mathbf{0}\vert D_0)\left[1-\Prob(\mathbf{0}\vert D_0)\right]^{T-k}\nonumber\\
&&~~~~~~~~~~~~\times\rme^{-\rmi h\,\hat{D}_{1} \,D_{0} (T-k)}.
%
\end{eqnarray}
%
Furthermore, using similar steps we obtain 
%
\begin{eqnarray}
&&\sum_{\Sm(\ell)}\Prob\left[\Sm(\ell)\vert D_\ell\right]\,\rme^{-\rmi T\hat{D}_{\ell+1} D_{\ell+1}[\Sm(\ell)]}\nonumber\\
&&~=\rme^{-\rmi T\hat{D}_{\ell+1} \left\{D_{\ell}-h\,D_{\ell}f \right\}}\sum_{k=0}^T{T\choose k}\Prob^k(\mathbf{0}\vert D_\ell)\left[1-\Prob(\mathbf{0}\vert D_\ell)\right]^{T-k}\nonumber\\
&&~~~~~~~~~~~~\times\rme^{-\rmi h\,\hat{D}_{\ell+1} \,D_{\ell} (T-k)},
%
\end{eqnarray}
%
where in above we defined the prob. 
$\Prob(\mathbf{0}\vert D_\ell)=\prod_{i=1}^N\left[1-\phi(w_i/D_\ell)\right]$.

%
\begin{eqnarray}
&&\prod_{\ell=0}^{L-1}\sum_{\Sm(\ell)}\Prob\left[\Sm(\ell)\vert D_\ell\right]\,\rme^{-\rmi T\hat{D}_{\ell+1} D_{\ell+1}[\Sm(\ell)]}\nonumber\\
&&~=\prod_{\ell=0}^{L-1}\rme^{-\rmi T\hat{D}_{\ell+1} \left\{D_{\ell}-h\,D_{\ell}f \right\}}\sum_{k=0}^T{T\choose k}\Prob^k(\mathbf{0}\vert D_\ell)\left[1-\Prob(\mathbf{0}\vert D_\ell)\right]^{T-k}\nonumber\\
&&~~~~~~~~~~~~\times\rme^{-\rmi h\,\hat{D}_{\ell+1} \,D_{\ell} (T-k)}.
%
\end{eqnarray}
%
Now consider 
%
\begin{eqnarray}
&&\Prob\left[D_{L},\ldots,D_{1}\vert D_{0}\right]\nonumber\\
%
&&=\left\{\prod_{\ell=1}^L\int_{-\infty}^{\infty}\frac{\rmd\hat{D}_\ell}{2\pi/T}\rme^{\rmi T\hat{D}_\ell D_{\ell}}\right\}\nonumber\\
&&\times\prod_{\ell=0}^{L-1}\sum_{k=0}^T{T\choose k}\Prob^k(\mathbf{0}\vert D_\ell)\left[1-\Prob(\mathbf{0}\vert D_\ell)\right]^{T-k}\nonumber\\
&&~~~\times\rme^{-\rmi T\hat{D}_{\ell+1} \left\{D_{\ell}-h\,D_{\ell}f \right\}-\rmi h\,\hat{D}_{\ell+1} \,D_{\ell} (T-k)}\nonumber\\
%
&&=\left\{\prod_{\ell=1}^L\int_{-\infty}^{\infty}\frac{\rmd\hat{D}_\ell}{2\pi/T}\rme^{\rmi T\hat{D}_\ell D_{\ell}}\right\}\prod_{\ell=0}^{L-1}\sum_{k_\ell=0}^T\Prob[k_\ell\vert D_\ell]\nonumber\\
&&~~~\times\prod_{\ell=0}^{L-1}\rme^{-\rmi T\hat{D}_{\ell+1} \left\{D_{\ell}-h\,D_{\ell}f \right\}-\rmi h\,\hat{D}_{\ell+1} \,D_{\ell} (T-k_\ell)}\nonumber\\
%
&&=\left\{\prod_{\ell=1}^L\int_{-\infty}^{\infty}\frac{\rmd\hat{D}_\ell}{2\pi/T}\rme^{\rmi T\hat{D}_\ell D_{\ell}}\right\}\prod_{\ell=0}^{L-1}\sum_{k_\ell=0}^T\Prob[k_\ell\vert D_\ell]\nonumber\\
&&~~~\times\prod_{\ell=1}^{L}\rme^{-\rmi T\hat{D}_{\ell} \left\{D_{\ell-1}-h\,D_{\ell-1}f \right\}-\rmi h\,\hat{D}_{\ell} \,D_{\ell-1} (T-k_{\ell-1})}\nonumber\\
%
&&+\prod_{\ell=0}^{L-1}\sum_{k_\ell=0}^T\Prob[k_\ell\vert D_\ell]\nonumber\\
&&~~~\times\prod_{\ell=1}^{L}\int_{-\infty}^{\infty}\frac{\rmd\hat{D}_\ell}{2\pi/T}\rme^{\rmi T\hat{D}_\ell D_{\ell}-\rmi T\hat{D}_{\ell} \left\{D_{\ell-1}-h\,D_{\ell-1}f \right\}-\rmi T\hat{D}_{\ell} \,h\,D_{\ell-1} (1-k_{\ell-1}/T)}\nonumber\\
%
&&=\prod_{\ell=0}^{L-1}\sum_{k_\ell=0}^T\Prob[k_\ell\vert D_\ell]\nonumber\\
&&~~~\times\prod_{\ell=1}^{L}\int_{-\infty}^{\infty}\frac{\rmd\hat{D}_\ell}{2\pi/T}\rme^{\rmi T\hat{D}_\ell \left[D_{\ell}- \left\{D_{\ell-1}-h\,D_{\ell-1}f \right\}- \,h\,D_{\ell-1} (1-k_{\ell-1}/T)\right]}\nonumber\\
%
&&=\prod_{\ell=0}^{L-1}\sum_{k_\ell=0}^T\Prob[k_\ell\vert D_\ell]\nonumber\\
&&~~~\times\prod_{\ell=1}^{L}\delta\left(D_{\ell}- \left\{D_{\ell-1}-h\,D_{\ell-1}f \right\}- \,h\,D_{\ell-1} (1-k_{\ell-1}/T)\right)\nonumber\\
%
&&=\prod_{\ell=0}^{L-1}\sum_{k_\ell=0}^T\Prob[k_\ell\vert D_\ell]\nonumber\\
&&~~~\times\prod_{\ell=0}^{L-1}\delta\left(D_{\ell+1}- \left\{D_{\ell}-h\,D_{\ell}f \right\}- \,h\,D_{\ell} (1-k_{\ell}/T)\right),
%
\end{eqnarray}
%
where in above we defined 
%
\begin{eqnarray}
\Prob[k\vert D_\ell]={T\choose k}\left[1-\phi\left(\sum_{i=1}^Nw_i/D_\ell\right)\right]^k\phi^{T-k}\left(\sum_{i=1}^Nw_i/D_\ell\right).
\end{eqnarray}
%
Thus from above follows that 
%
\begin{eqnarray}
\Prob\left[D_{L},\ldots,D_{1}\vert D_{0}\right]
%
&=&\prod_{\ell=0}^{L-1}\sum_{k_\ell=0}^T\Prob[k_\ell\vert D_\ell]\nonumber\\
&&\times\prod_{\ell=0}^{L-1}\delta\left(D_{\ell+1}- D_{\ell}+ h\,D_{\ell}\left(f- 1+\frac{k_{\ell}}{T}\right)\right),
%
\end{eqnarray}
%
where $\Prob[k\vert D]$ is the binomial distribution
%
\begin{eqnarray}
\Prob[k\vert D]={T\choose k}\left[1-\phi\left(\sum_{i=1}^Nw_i/D\right)\right]^k\phi^{T-k}\left(\sum_{i=1}^Nw_i/D\right).
\end{eqnarray}
%
\end{proof}



\subsection{Saddle-point mean-field theory}
%
For a test function $\mathcal{F}[D_{L},\ldots,D_{1}]$ we consider the average 
%
\begin{eqnarray}
%
&&\langle\mathcal{F}[D_{L},\ldots,D_{1}]\rangle\nonumber\\
&&~~~~~~~~~~~~~~=\int\rmd D_{L}\cdots\int\rmd D_{1}\Prob\left[D_{L},\ldots,D_{1}\vert D_{0}\right]
\,\mathcal{F}[D_{L},\ldots,D_{1}]\label{def:F-aver}
%
\end{eqnarray}
%
as follows
%
\begin{eqnarray}
&&\langle\mathcal{F}[D_{L},\ldots,D_{1}]\rangle\nonumber\\
%
&=&\int\rmd D_{L}\cdots\int\rmd D_{1}\prod_{\ell=0}^{L-1}\sum_{k_\ell=0}^T\Prob[k_\ell\vert D_\ell]\nonumber\\
&&\times\prod_{\ell=0}^{L-1}\delta\left(D_{\ell+1}- D_{\ell}+ h\,D_{\ell}\left(f- 1+\frac{k_{\ell}}{T}\right)\right)\nonumber\\
%
&&\times \mathcal{F}[D_{L},\ldots,D_{1}]\nonumber\\
%
&=&\int\rmd D_{L}\cdots\int\rmd D_{1}\prod_{\ell=0}^{L-1}\sum_{k_\ell=0}^T\Prob[k_\ell\vert D_\ell]\nonumber\\
&&\times\prod_{\ell=0}^{L-1}\int_{-\infty}^{\infty}\frac{\rmd\hat{D}_\ell}{2\pi/T}\rme^{\rmi T\hat{D}_\ell\left(D_{\ell+1}- D_{\ell}+ h\,D_{\ell}\left(f- 1+\frac{k_{\ell}}{T}\right)\right)}
\nonumber\\
%
&&\times \mathcal{F}[D_{L},\ldots,D_{1}]\nonumber\\
%
&=&\int\rmd D_{L}\cdots\int\rmd D_{1}\prod_{\ell=0}^{L-1}\int_{-\infty}^{\infty}\frac{\rmd\hat{D}_\ell}{2\pi/T}\rme^{\rmi T\hat{D}_\ell\left(D_{\ell+1}- D_{\ell}+ h\,D_{\ell}\left(f- 1\right)\right)}\nonumber\\
&&\times\prod_{\ell=0}^{L-1}\sum_{k_\ell=0}^T\Prob[k_\ell\vert D_\ell] \rme^{\rmi \hat{D}_\ell h\,D_{\ell}k_{\ell} }
\nonumber\\
%
&&\times \mathcal{F}[D_{L},\ldots,D_{1}]\nonumber\\
%
&=&\int\rmd D_{L}\cdots\int\rmd D_{1}\prod_{\ell=0}^{L-1}\int_{-\infty}^{\infty}\frac{\rmd\hat{D}_\ell}{2\pi/T}\rme^{\rmi T\hat{D}_\ell\left(D_{\ell+1}- D_{\ell}+ h\,D_{\ell}\left(f- 1\right)\right)}\nonumber\\
&&\times\prod_{\ell=0}^{L-1}\left\{ \phi\left(\sum_{i=1}^Nw_i/D_\ell\right)+\left[1-\phi\left(\sum_{i=1}^Nw_i/D_\ell\right)\right]\rme^{\rmi \hat{D}_\ell h\,D_{\ell}}\right\}^T
\nonumber\\
%
&&\times \mathcal{F}[D_{L},\ldots,D_{1}]\nonumber\\
%
&=&\int\rmd D_{L}\cdots\int\rmd D_{1}\prod_{\ell=0}^{L-1}\int_{-\infty}^{\infty}\frac{\rmd\hat{D}_\ell}{2\pi/T}\nonumber\\
&&\times\rme^{T\Psi\left[D_{L},\ldots,D_{1};\hat{D}_{L-1},\ldots,\hat{D}_0\right]} \mathcal{F}[D_{L},\ldots,D_{1}]
%
\end{eqnarray}
%
where in above we defined 
%
\begin{eqnarray}
&&\Psi\left[D_{L},\ldots,D_{1};\hat{D}_{L-1},\ldots,\hat{D}_0\right]\nonumber\\
%
&&~~~=\rmi \sum_{\ell=0}^{L-1}\hat{D}_\ell\left(D_{\ell+1}- D_{\ell}+ h\,D_{\ell}\left(f- 1\right)\right)\nonumber\\
&&~~~~~+\sum_{\ell=0}^{L-1}\log\left( \phi\left(\sum_{i=1}^Nw_i/D_\ell\right)+\left[1-\phi\left(\sum_{i=1}^Nw_i/D_\ell\right)\right]\rme^{\rmi \hat{D}_\ell h\,D_{\ell}}\right)\label{def:Psi}.
%
\end{eqnarray}
%
Now using the normalisation property 
%
\begin{eqnarray}
%
&&\int\rmd D_{L}\cdots\int\rmd D_{1}\Prob\left[D_{L},\ldots,D_{1}\vert D_{0}\right]\nonumber\\
%
&&=\int\rmd D_{L}\cdots\int\rmd D_{1}\prod_{\ell=0}^{L-1}\int_{-\infty}^{\infty}\frac{\rmd\hat{D}_\ell}{2\pi/T}\rme^{T\Psi\left[D_{L},\ldots,D_{1};\hat{D}_{L-1},\ldots,\hat{D}_0\right]} =1
%
\end{eqnarray}
%
we obtain
%
\begin{eqnarray}
&&\langle\mathcal{F}[D_{L},\ldots,D_{1}]\rangle\nonumber\\
%
&&~~~~~~~~~~~~=\mathcal{C}_T\int\rmd D_{L}\cdots\int\rmd D_{1}\int\rmd \hat{D}_{L-1}\cdots\int\rmd \hat{D}_{0}\nonumber\\
 &&~~~~~~~~~~~~~~~~~\times\rme^{T\Psi\left[D_{L},\ldots,D_{1};\hat{D}_{L-1},\ldots,\hat{D}_0\right]}\mathcal{F}[D_{L},\ldots,D_{1}]\label{eq:SP-integral},
%
\end{eqnarray}
%
where we defined
%
\begin{eqnarray}
%
\mathcal{C}_T&=&\int\!\rmd D_{L}\!\cdots\!\int\!\rmd D_{1}\int\!\rmd \hat{D}_{L-1}\!\cdots\!\int\!\rmd \hat{D}_{0} \rme^{T\Psi\left[D_{L},\ldots,D_{1};\hat{D}_{L-1},\ldots,\hat{D}_0\right]}.
%
\end{eqnarray}
%
For $T\to\infty$ the integral (\ref{eq:SP-integral}) is dominated by extremum of $\Psi\left[\cdots;\cdots\right]$. To find the latter we consider the derivative
%
\begin{eqnarray}
&&\nonumber\\
%
&&\frac{\partial}{\partial \hat{D}_\ell}\Psi\left[\cdots;\cdots\right]=\rmi \left(D_{\ell+1}- D_{\ell}+ h\,D_{\ell}\left(f- 1\right)\right)\nonumber\\
&&~~~~~~~~~~~~~+ \frac{\rmi h\,D_{\ell}\left[1-\phi\left(\sum_{i=1}^Nw_i/D_\ell\right)\right]\rme^{\rmi \hat{D}_\ell h\,D_{\ell}}}{\phi\left(\sum_{i=1}^Nw_i/D_\ell\right)+\left[1-\phi\left(\sum_{i=1}^Nw_i/D_\ell\right)\right]\rme^{\rmi \hat{D}_\ell h\,D_{\ell}}}\label{eq:Psi-der-D-hat}.
%
\end{eqnarray}
%
then solving $\frac{\partial}{\partial \hat{D}_\ell}\Psi\left[\cdots;\cdots\right]=0$ gives us the equation 
%
\begin{eqnarray}
%
&&D_{\ell+1}=D_{\ell}- h\,D_{\ell}\left(f- 1\right)\nonumber\\
&&~~~~~~~~~~~~~- \frac{h\,D_{\ell}\left[1-\phi\left(\sum_{i=1}^Nw_i/D_\ell\right)\right]\rme^{\rmi \hat{D}_\ell h\,D_{\ell}}}{\phi\left(\sum_{i=1}^Nw_i/D_\ell\right)+\left[1-\phi\left(\sum_{i=1}^Nw_i/D_\ell\right)\right]\rme^{\rmi \hat{D}_\ell h\,D_{\ell}}}\label{eq:SP-D}.
%
\end{eqnarray}
%
Furthermore, $\frac{\partial}{\partial D_L}\Psi\left[\cdots;\cdots\right]=\rmi \hat{D}_{L-1}$ and hence solving $\frac{\partial}{\partial D_L}\Psi\left[\cdots;\cdots\right]=0$ gives us $ \hat{D}_{L-1}=0$. Using the latter in (\ref{eq:SP-D}) we obtain the equation 
%
\begin{eqnarray}
%
D_{L}&=&D_{L-1}- h\,D_{L-1}\left(f- \phi\left(\sum_{i=1}^Nw_i/D_{L-1}\right)\right)\label{eq:SP-D-L+1}.
%
\end{eqnarray}
%
Let us consider the derivative
%
\begin{eqnarray}
&&\frac{\partial}{\partial D_{L-1}}\Psi\left[D_{L},\ldots,D_{1};\hat{D}_{L-1},\ldots,\hat{D}_0\right]\nonumber\\
%
&&~~~~~~~~~~~~~~=\rmi \hat{D}_{L-1}\left(- 1+ h\,\left(f- 1\right)\right)+\rmi \hat{D}_{L-2}\nonumber\\
&&~~~~~~~~~~~~~~~~~~~~~~+\frac{\partial}{\partial D_{L-1}}\log\left( \phi+\left[1-\phi\right]\rme^{\rmi \hat{D}_{L-1} h\,D_{L-1} }\right),
%
\end{eqnarray}
%
where in above $\phi\equiv\phi\left(\sum_{i=1}^Nw_i/D_{L-1}\right)$. The second term in above can be computed as follows
%
\begin{eqnarray}
%
&&\frac{\partial}{\partial D_{L-1}}\log\left( \phi+\left[1-\phi\right]\rme^{\rmi \hat{D}_{L-1} h\,D_{L-1} }\right)\nonumber\\
%
&&=
\frac{\frac{\partial}{\partial D_{L-1}}\phi+\frac{\partial}{\partial D_{L-1}}\left[1-\phi\right]\rme^{\rmi \hat{D}_{L-1} h\,D_{L-1} }}{\phi+\left[1-\phi\right]\rme^{\rmi \hat{D}_{L-1} h\,D_{L-1} }}\nonumber\\
%
&&=
\frac{\frac{\partial}{\partial D_{L-1}}\phi\left[1-\rme^{\rmi \hat{D}_{L-1} h\,D_{L-1} }\right]+\left[1-\phi\right]\rme^{\rmi \hat{D}_{L-1} h\,D_{L-1} }\rmi \hat{D}_{L-1} h}{\phi+\left[1-\phi\right]\rme^{\rmi \hat{D}_{L-1} h\,D_{L-1} }}.
%
\end{eqnarray}
%
The condition for extremum $\frac{\partial}{\partial D_{L-1}}\Psi\left[D_{L},\ldots,D_{1};\hat{D}_{L-1},\ldots,\hat{D}_0\right]=0$ gives us the equation
%
\begin{eqnarray}
&&\rmi \hat{D}_{L-1}\left(- 1+ h\,\left(f- 1\right)\right)+\rmi \hat{D}_{L-2}\nonumber\\
&&~~~=-\frac{\frac{\partial}{\partial D_{L-1}}\phi\left[1-\rme^{\rmi \hat{D}_{L-1} h\,D_{L-1} }\right]+\left[1-\phi\right]\rme^{\rmi \hat{D}_{L-1} h\,D_{L-1} }\rmi \hat{D}_{L-1} h}{\phi+\left[1-\phi\right]\rme^{\rmi \hat{D}_{L-1} h\,D_{L-1} }}.
%
\end{eqnarray}
%
However, we found earlier that $\hat{D}_{L-1}=0$ and hence above gives us $\hat{D}_{L-2}=0$. The latter implies that $\hat{D}_{\ell}=0$ for $\ell\in\{0,1,\ldots,L-1\}$ and using this in the equation (\ref{eq:SP-D-simpl}) gives us 
%
\begin{eqnarray}
%
D_{\ell+1}&=&D_{\ell}- h\,D_{\ell}\left(f- \phi\left(\frac{\sum_{i=1}^Nw_i}{D_{\ell}}\right)\right)\label{eq:SP-D-simpl}.
%
\end{eqnarray}
%
We note that on one hand, by (\ref{eq:SP-integral}), we have 
%
\begin{eqnarray}
%
\lim_{T\to\infty}
\langle\mathcal{F}[D_{L},\ldots,D_{1}]\rangle&=&\mathcal{F}[D_{L},\ldots,D_{1}],
%
\end{eqnarray}
%
where $D_{\ell}$ variables on the right hand side are computed by the equation (\ref{eq:SP-D-simpl}), but on other hand for  $\mathcal{F}[D_{L},\ldots,D_{1}]=D_{\ell}$ one the left hand side we have $\langle D_{\ell} \rangle=D_{\ell}$. Hence from the latter follows that  (\ref{eq:SP-D-simpl}) is the mean-field equation 
%
%
\begin{eqnarray}
%
\langle D_{\ell+1}\rangle&=&\langle D_{\ell}\rangle- h\,\langle D_{\ell}\rangle\left(f- \phi\left(\frac{\sum_{i=1}^Nw_i}{\langle D_{\ell}\rangle}\right)\right)\label{eq:MF}.
%
\end{eqnarray}
%


\subsection{Mean-field approximation and concentration}
%
Let
\begin{equation}
W=\sum_{i=1}^N w_i,
\qquad
\phi(u)=1-(1-f)^u,
\qquad
0<f<1.
\label{eq:phi_def}
\end{equation}

We consider the stochastic recurrence
\begin{equation}
D_{\ell+1}
=
D_\ell-hD_\ell
\left(
f-1+\frac{k_\ell}{T}
\right),
\label{eq:stochastic_recursion}
\end{equation}
where, conditionally on $D_\ell$,
\begin{equation}
k_\ell\mid D_\ell
\sim
\mathrm{Bin}(T,p(D_\ell)),
\label{eq:binomial_model}
\end{equation}
with
\begin{equation}
p(D)=1-\phi(W/D).
\label{eq:p_def}
\end{equation}

We define the deterministic mean-field map
\begin{equation}
g(D)
=
D-hD(f-\phi(W/D)),
\label{eq:g_def}
\end{equation}
and the deterministic mean-field recursion
\begin{equation}
d_{\ell+1}=g(d_\ell),
\qquad
d_0=D_0>0.
\label{eq:mean_field_recursion}
\end{equation}

Assume
\begin{equation}
h<\frac1f.
\label{eq:h_condition}
\end{equation}

For a fixed finite horizon $L\ge1$, define
\begin{equation}
D_{\min}
:=
D_0(1-hf)^L,
\label{eq:Dmin_def}
\end{equation}
and
\begin{equation}
D_{\max}
:=
D_0(1+h(1-f))^L.
\label{eq:Dmax_def}
\end{equation}

Let
\begin{equation}
A=-\log(1-f)>0.
\label{eq:c_def}
\end{equation}

We define the critical value $D_c$ by the condition
\begin{equation}
g'(D_c)=1.
\label{eq:Dc_definition}
\end{equation}

Solving this equation yields
\begin{equation}
D_c
=
\frac{
W\,A
}{
-\mathcal W_{-1}\!\left(-\frac{1-f}{e}\right)-1
},
\label{eq:Dc_formula}
\end{equation}
where $\mathcal W_{-1}$ denotes the $-1$ branch of the Lambert $W$-function.

We assume
\begin{equation}
D_{\min}>D_c.
\label{eq:contractive_assumption}
\end{equation}

\begin{theorem}
\label{th:concentration}
Assume (\ref{eq:h_condition}) and (\ref{eq:contractive_assumption}). Then, for every fixed horizon $L$,
\begin{equation}
\sup_{0\le \ell\le L}
|D_\ell-d_\ell|
=
O_{\mathbb P}(T^{-1/2}).
\label{eq:main_concentration}
\end{equation}

Furthermore, if
\begin{equation}
m_\ell:=\mathbb E[D_\ell],
\label{eq:mell_def}
\end{equation}
then
\begin{equation}
m_\ell-d_\ell
=
O(T^{-1/2}),
\qquad
0\le \ell\le L.
\label{eq:mean_estimate}
\end{equation}
\end{theorem}

\begin{proof}

Since
\begin{equation}
0\le \frac{k_\ell}{T}\le 1,
\label{eq:k_bounds}
\end{equation}
we obtain
\begin{equation}
f-1
\le
f-1+\frac{k_\ell}{T}
\le
f.
\label{eq:middle_bounds}
\end{equation}

Multiplying (\ref{eq:middle_bounds}) by $-h$ gives
\begin{equation}
-hf
\le
-h\left(f-1+\frac{k_\ell}{T}\right)
\le
h(1-f).
\label{eq:middle_bounds2}
\end{equation}

Adding $1$ to (\ref{eq:middle_bounds2}) yields
\begin{equation}
1-hf
\le
1-h\left(f-1+\frac{k_\ell}{T}\right)
\le
1+h(1-f).
\label{eq:middle_bounds3}
\end{equation}

Using (\ref{eq:stochastic_recursion}), we conclude that
\begin{equation}
D_\ell(1-hf)
\le
D_{\ell+1}
\le
D_\ell(1+h(1-f)).
\label{eq:D_bounds_step}
\end{equation}

Iterating (\ref{eq:D_bounds_step}) gives
\begin{equation}
D_0(1-hf)^\ell
\le
D_\ell
\le
D_0(1+h(1-f))^\ell.
\label{eq:D_bounds_iterated}
\end{equation}

Since (\ref{eq:h_condition}) implies
\begin{equation}
1-hf>0,
\label{eq:positive_factor}
\end{equation}
we obtain from (\ref{eq:D_bounds_iterated}) that
\begin{equation}
D_{\min}\le D_\ell\le D_{\max},
\qquad
0\le \ell\le L.
\label{eq:D_interval}
\end{equation}

Similarly, the deterministic recursion (\ref{eq:mean_field_recursion}) satisfies
\begin{equation}
d_{\ell+1}
=
d_\ell
\left(
1-h(f-\phi(W/d_\ell))
\right).
\label{eq:d_step}
\end{equation}

Since
\begin{equation}
0\le \phi(W/d_\ell)\le 1,
\label{eq:phi_bounds}
\end{equation}
we obtain
\begin{equation}
1-hf
\le
1-h(f-\phi(W/d_\ell))
\le
1+h(1-f),
\label{eq:d_factor_bounds}
\end{equation}
and hence
\begin{equation}
D_{\min}\le d_\ell\le D_{\max},
\qquad
0\le \ell\le L.
\label{eq:d_interval}
\end{equation}

We now analyze the derivative of $g$. Differentiating (\ref{eq:g_def}) gives
\begin{equation}
g'(D)
=
1-hf+h\phi(W/D)
-h\frac{W}{D}\phi'(W/D).
\label{eq:gprime1}
\end{equation}

From (\ref{eq:phi_def}),
\begin{equation}
\phi'(u)
=
-(1-f)^u\log(1-f).
\label{eq:phi_prime}
\end{equation}

Substituting (\ref{eq:phi_prime}) into (\ref{eq:gprime1}) yields
\begin{equation}
g'(D)
=
1-hf+h\phi(W/D)
+
h\frac{W}{D}(1-f)^{W/D}\log(1-f).
\label{eq:gprime2}
\end{equation}

Let
\begin{equation}
u=\frac{W}{D}.
\label{eq:u_def}
\end{equation}

Using (\ref{eq:c_def}), we have
\begin{equation}
(1-f)^u=e^{-Au},
\label{eq:exp_form}
\end{equation}
and therefore
\begin{equation}
g'(D)
=
1-hf+hq(u),
\label{eq:gprime_q}
\end{equation}
where
\begin{equation}
q(u)=1-(1+Au)e^{-Au}.
\label{eq:q_def}
\end{equation}

Differentiating (\ref{eq:q_def}) gives
\begin{equation}
q'(u)=A^2ue^{-Au}>0.
\label{eq:qprime}
\end{equation}

Hence $q(u)$ is strictly increasing in $u$. Since $u=W/D$ is strictly decreasing in $D$, equation (\ref{eq:gprime_q}) implies that $g'(D)$ is strictly decreasing in $D$.

Moreover,
\begin{equation}
q(u)\ge0,
\label{eq:q_positive}
\end{equation}
and therefore
\begin{equation}
g'(D)\ge 1-hf>0.
\label{eq:gprime_positive}
\end{equation}

Thus
\begin{equation}
|g'(D)|=g'(D).
\label{eq:absolute_removed}
\end{equation}

The critical value $D_c$ is defined by (\ref{eq:Dc_definition}), namely
\begin{equation}
1-hf+hq(W/D_c)=1.
\label{eq:Dc_equation}
\end{equation}

Since $h>0$, this reduces to
\begin{equation}
q(W/D_c)=f.
\label{eq:Dc_equation2}
\end{equation}

Let
\begin{equation}
u_c=\frac{W}{D_c}.
\label{eq:uc_def}
\end{equation}

Using (\ref{eq:q_def}), equation (\ref{eq:Dc_equation2}) becomes
\begin{equation}
1-(1+Au_c)e^{-Au_c}=f.
\label{eq:uc_equation}
\end{equation}

Equivalently,
\begin{equation}
(1+Au_c)e^{-Au_c}=1-f.
\label{eq:uc_equation2}
\end{equation}

Define
\begin{equation}
z=1+Au_c.
\label{eq:z_def}
\end{equation}

Then (\ref{eq:uc_equation2}) becomes
\begin{equation}
ze^{-z}=\frac{1-f}{e}.
\label{eq:z_equation}
\end{equation}

Multiplying by $-1$ and applying the Lambert $W$-function gives
\begin{equation}
-z
=
\mathcal W_{-1}
\left(
-\frac{1-f}{e}
\right).
\label{eq:z_solution}
\end{equation}

Substituting back yields (\ref{eq:Dc_formula}).

Since $g'(D)$ is strictly decreasing and $g'(D_c)=1$, assumption (\ref{eq:contractive_assumption}) implies
\begin{equation}
0<g'(D)<1,
\qquad
D\in[D_{\min},D_{\max}].
\label{eq:contractive_region}
\end{equation}

Define
\begin{equation}
L_g
:=
\sup_{D\in[D_{\min},D_{\max}]}
|g'(D)|.
\label{eq:Lg_def}
\end{equation}

Using (\ref{eq:contractive_region}), we obtain
\begin{equation}
L_g<1.
\label{eq:Lg_less_1}
\end{equation}

Define the fluctuation
\begin{equation}
\epsilon_\ell=D_\ell-d_\ell.
\label{eq:error_def}
\end{equation}

Write
\begin{equation}
\frac{k_\ell}{T}
=
p(D_\ell)+\eta_\ell,
\label{eq:eta_decomposition}
\end{equation}
where
\begin{equation}
\eta_\ell
=
\frac{k_\ell}{T}-p(D_\ell).
\label{eq:eta_def}
\end{equation}

Substituting (\ref{eq:eta_decomposition}) into (\ref{eq:stochastic_recursion}) gives
\begin{equation}
D_{\ell+1}
=
g(D_\ell)-hD_\ell\eta_\ell.
\label{eq:D_with_eta}
\end{equation}

Subtracting (\ref{eq:mean_field_recursion}) from (\ref{eq:D_with_eta}), we obtain
\begin{equation}
\epsilon_{\ell+1}
=
g(D_\ell)-g(d_\ell)
-
hD_\ell\eta_\ell.
\label{eq:error_recursion}
\end{equation}

Using (\ref{eq:Lg_less_1}),
%
\begin{equation}
|g(D_\ell)-g(d_\ell)|
\le
L_g|\epsilon_\ell|.
\label{eq:lipschitz_estimate}
\end{equation}

Combining (\ref{eq:error_recursion}), (\ref{eq:lipschitz_estimate}), and (\ref{eq:D_interval}), we obtain
\begin{equation}
|\epsilon_{\ell+1}|
\le
L_g|\epsilon_\ell|
+
hD_{\max}|\eta_\ell|.
\label{eq:error_recursive_bound}
\end{equation}

\begin{comment}
\begin{eqnarray}
|e_{1}|
&\le&
hD_{\max}|\eta_0|\\
%
|e_{2}|
&\le&
L_g|e_1|
+
hD_{\max}|\eta_1| \le hD_{\max} (L_g|\eta_0|
+
|\eta_1| )\\
%
|e_{3}|
&\le&
L_g|e_2|
+
hD_{\max}|\eta_2|\le  hD_{\max} (L^2_g|\eta_0|
+
L_g|\eta_1| 
+
|\eta_2|)
%
\end{eqnarray}
\end{comment}

Since $\epsilon_0=0$,  iterating (\ref{eq:error_recursive_bound}) yields
%
\begin{equation}
|\epsilon_\ell|
\le
hD_{\max}
\sum_{r=0}^{\ell-1}
L_g^{\ell-1-r}
|\eta_r|.
\label{eq:error_series}
\end{equation}

From the binomial model (\ref{eq:binomial_model}),
\begin{equation}
\mathbb E[\eta_\ell\mid D_\ell]=0,
\label{eq:eta_mean}
\end{equation}
and
\begin{equation}
\mathrm{Var}(\eta_\ell\mid D_\ell)
=
\frac{p(D_\ell)(1-p(D_\ell))}{T}.
\label{eq:eta_variance}
\end{equation}

Since $p(1-p)\le \frac14$, we obtain
\begin{equation}
\mathrm{Var}(\eta_\ell\mid D_\ell)
\le
\frac1{4T}.
\label{eq:eta_variance_bound}
\end{equation}
%
Also, $|\eta_\ell|\le1$.


Let us define
%
\begin{equation}
a_{\ell,r}
=
hD_{\max}L_g^{\ell-1-r}
\label{eq:a_def}
\end{equation}
%
then (\ref{eq:error_series}) becomes
\begin{equation}
|\epsilon_\ell|
\le
\sum_{r=0}^{\ell-1}
a_{\ell,r}|\eta_r|.
\label{eq:error_series2}
\end{equation}
%
Squaring both sides in above and applying the Cauchy--Schwarz inequality  gives us
%
\begin{equation}
|\epsilon_\ell|^2
\le
\left(
\sum_{r=0}^{\ell-1}a_{\ell,r}^2
\right)
\left(
\sum_{r=0}^{\ell-1}\eta_r^2
\right).
\label{eq:cauchy_estimate}
\end{equation}

Since $L_g<1$, we have 
\begin{equation}
\sum_{r=0}^{\ell-1}a_{\ell,r}^2
=
h^2D_{\max}^2
\sum_{r=0}^{\ell-1}
L_g^{2(\ell-1-r)}
\le
\frac{h^2D_{\max}^2}{1-L_g^2}.
\label{eq:geometric_bound}
\end{equation}

Taking expectation in (\ref{eq:cauchy_estimate}), using (\ref{eq:eta_variance_bound}), and recalling that $L$ is fixed, we obtain
\begin{equation}
\mathbb E[\epsilon_\ell^2]
=
O(T^{-1}).
\label{eq:e_second_moment}
\end{equation}

By Chebyshev's inequality,
\begin{equation}
\vert\epsilon_\ell\vert
=
O_{\mathbb P}(T^{-1/2}).
\label{eq:e_op}
\end{equation}

Since $L$ is fixed,
\begin{equation}
\sup_{0\le \ell\le L}|\epsilon_\ell|
=
O_{\mathbb P}(T^{-1/2}).
\label{eq:sup_e}
\end{equation}

Using (\ref{eq:error_def}), we obtain (\ref{eq:main_concentration}).

Finally,
\begin{equation}
m_\ell-d_\ell
=
\mathbb E[\epsilon_\ell].
\label{eq:m_minus_d}
\end{equation}

Applying Cauchy--Schwarz and using (\ref{eq:e_second_moment}),
\begin{equation}
|m_\ell-d_\ell|
\le
\mathbb E|\epsilon_\ell|
\le
\sqrt{\mathbb E[\epsilon_\ell^2]}
=
O(T^{-1/2}),
\label{eq:mean_bound_final}
\end{equation}
which proves (\ref{eq:mean_estimate}).
\end{proof}

Theorem~\ref{th:concentration} shows that, on any fixed time horizon, once the estimator remains above the critical threshold $D_c$, the stochastic recurrence (\ref{eq:stochastic_recursion}) is well approximated by the deterministic mean-field recursion (\ref{eq:mean_field_recursion}), with fluctuations of order $T^{-1/2}$. Thus the saddle-point mean-field equation is replaced by a rigorous concentration result. The condition $D_{\min}>D_c$ has a clear dynamical interpretation: the estimator remains inside the contractive region of the map $g$, where stochastic perturbations are damped rather than amplified. Moreover, the threshold becomes particularly mild in the sparse-block regime. Indeed, from (\ref{eq:Dc_formula}),
\begin{equation}
\frac{D_c}{W}
=
\frac{
\log(1-f)
}{
\mathcal W_{-1}\!\left(-\frac{1-f}{e}\right)+1
},
\label{eq:Dc_over_W}
\end{equation}
one obtains the asymptotic behaviour
\begin{equation}
\frac{D_c}{W}
\sim
\sqrt{\frac{f}{2}},
\qquad
f\to0.
\label{eq:Dc_asymptotic}
\end{equation}
Hence
\begin{equation}
\frac{D_c}{W}\to0
\qquad
\mbox{as}
\qquad
f\to0,
\label{eq:Dc_limit}
\end{equation}
which implies that, for small active-slot coefficient $f$, the contractive region $D>D_c$ occupies almost the entire positive half-line. In particular, for the practically relevant value $f=1/30$, equation (\ref{eq:Dc_formula}) gives
\begin{equation}
\frac{D_c}{W}\approx 0.121,
\label{eq:Dc_numerical}
\end{equation}
showing that the estimator may substantially underestimate the true stake and still remain in the stable concentration regime.


\section*{Acknowledgements}
Authors would like to thank... 

\appendix

\section{Average}
%
Let us consider
%
\begin{eqnarray}
&&\int\rmd D_{L}\cdots\int\rmd D_{1}\Prob\left[D_{L},\ldots,D_{1}\vert D_{0}\right]
\,D_{L}\nonumber\\
&&=\int\rmd D_{L}\cdots\int\rmd D_{1}\prod_{\ell=0}^{L-1}\sum_{k_\ell=0}^T\Prob[k_\ell\vert D_\ell]\nonumber\\
&&\times\prod_{\ell=0}^{L-1}\delta\left(D_{\ell+1}- D_{\ell}+ h\,D_{\ell}\left[f- 1+\frac{k_{\ell}}{T}\right]\right)\nonumber\\
&&\times D_{L}\nonumber\\
%
&&=\int\rmd D_{L-1}\cdots\int\rmd D_{1}\prod_{\ell=0}^{L-1}\sum_{k_\ell=0}^T\Prob[k_\ell\vert D_\ell]\nonumber\\
&&\times\prod_{\ell=0}^{L-2}\delta\left(D_{\ell+1}- D_{\ell}+ h\,D_{\ell}\left[f- 1+\frac{k_{\ell}}{T}\right]\right)\nonumber\\
%
&&\times \left(D_{L-1}- h\,D_{L-1}\left[f- 1+\frac{k_{L-1}}{T}\right]\right)\nonumber\\
%
\end{eqnarray}
%
Let us define the function  

\begin{eqnarray}
%
F[D_{\ell},k_{\ell}]= D_{\ell}\left\{1- h\left[f- 1+\frac{k_{\ell}}{T}\right]\right\}\label{def:F}
%
\end{eqnarray}
%
and consider
%
\begin{eqnarray}
\Prob\left[k_{L-1},\ldots,k_{0}\vert D_{0}\right]
%
&=&\int\rmd D_{L}\cdots\int\rmd D_{1}\prod_{\ell=0}^{L-1}\Prob[k_\ell\vert D_\ell]\nonumber\\
&&\times\prod_{\ell=0}^{L-1}\delta\left(D_{\ell+1}- F[D_{\ell},k_{\ell}]\right)\nonumber\\
%
&=&\int\rmd D_{L-1}\cdots\int\rmd D_{1}\prod_{\ell=0}^{L-1}\Prob[k_\ell\vert D_\ell]\nonumber\\
&&\times\prod_{\ell=1}^{L-1}\delta\left(D_{\ell}- F[D_{\ell-1},k_{\ell-1}]\right)
%
\end{eqnarray}
%
For $L=2$ above gives us 
%
\begin{eqnarray}
\Prob\left[k_{1},k_{0}\vert D_{0}\right]
%
&=&\int\rmd D_{1}\prod_{\ell=0}^{1}\Prob[k_\ell\vert D_\ell]\delta\left(D_{1}- F[D_{0},k_{0}]\right)\nonumber\\
%
&=&\Prob[k_1\vert F[D_{0},k_{0}]]\,\Prob[k_0\vert D_0]. 
%
\end{eqnarray}
%
and for $L=3$ we get 
%
\begin{eqnarray}
\Prob\left[k_{2},k_{1},k_{0}\vert D_{0}\right]
%
&=&\int\rmd D_{2}\int\rmd D_{1}\prod_{\ell=0}^{2}\Prob[k_\ell\vert D_\ell]\nonumber\\
&&\times\prod_{\ell=1}^{2}\delta\left(D_{\ell}- F[D_{\ell-1},k_{\ell-1}]\right)\nonumber\\
%
&=&\int\rmd D_{2}\int\rmd D_{1}\Prob[k_2\vert D_2]\,\Prob[k_1\vert D_1]\,\Prob[k_0\vert D_0]\nonumber\\
&&\times\delta\left(D_{2}- F[D_{1},k_{1}]\right)\delta\left(D_{1} - F[D_{0},k_{0}]\right)\nonumber\\
%
&=&\Prob[k_2\vert F[F[D_{0},k_{0}],k_{1}]]\,\Prob[k_1\vert F[D_{0},k_{0}]]\,\Prob[k_0\vert D_0]
%
\end{eqnarray}
%
Consider 
%
\begin{eqnarray}
%
F[F[D_{0},k_{0}],k_{1}]&=& F[D_{0},k_{0}]\left\{1- h\,\left[f- 1+\frac{k_{1}}{T}\right]\right\}\nonumber\\
%
&=& D_{0}\left\{1- h\left[f- 1+\frac{k_{0}}{T}\right]\right\}\left\{1- h\,\left[f- 1+\frac{k_{1}}{T}\right]\right\}\nonumber\\
%
&=& D_{0}\prod_{\ell=0}^1\left\{1- h\left[f- 1+\frac{k_{\ell}}{T}\right]\right\}.
%
\end{eqnarray}
%
Using above
%
\begin{eqnarray}
%
F[F[F[D_{0},k_{0}],k_{1}],k_{2}]&=& F\left[D_{0}\prod_{\ell=0}^1\left\{1- h\left[f- 1+\frac{k_{\ell}}{T}\right]\right\},k_2\right]\nonumber\\
%
&=& D_{0}\prod_{\ell=0}^2\left\{1- h\left[f- 1+\frac{k_{\ell}}{T}\right]\right\}\nonumber\\
%
\end{eqnarray}
%
Now from above it is clear that 
%
\begin{eqnarray}
%
F[\cdots,k_{L}]&=& D_{0}\prod_{\ell=0}^L\left\{1- h\left[f- 1+\frac{k_{\ell}}{T}\right]\right\},
%
\end{eqnarray}
%
where  $F[\cdots,k_{0}]= D_{0}$, and from above follows
%
\begin{eqnarray}
\Prob\left[k_{L-1},\ldots,k_{0}\vert D_{0}\right]&=&\prod_{\ell=0}^{L-1}\Prob[k_\ell\vert F[\cdots,k_{\ell-1}]]\nonumber\\
%
&=&\prod_{\ell=0}^{L-1}\Prob\left[k_\ell\Bigg\vert D_{0}\prod_{\tilde{\ell}=0}^{\ell-1}\left\{1- h\left[f- 1+\frac{k_{\tilde{\ell}}}{T}\right]\right\}\right].
%
\end{eqnarray}
%
Finally, using the definition (\ref{def:F}) we consider
%
\begin{eqnarray}
\Prob\left[D_{L},\ldots,D_{1}\vert D_{0}\right]
%
&=&\prod_{\ell=0}^{L-1}\sum_{k_\ell=0}^T\Prob[k_\ell\vert D_\ell]\nonumber\\
&&\times\prod_{\ell=0}^{L-1}\delta\left(D_{\ell+1}- F[D_{\ell},k_{\ell}]\right)\nonumber\\
%
&=&\sum_{k_0=0}^T\cdots \sum_{k_{L-1}=0}^T\Prob\left[k_{L-1},\ldots,k_{0}\vert D_{0}\right]\nonumber\\
&&\times\prod_{\ell=0}^{L-1}\delta\left(D_{\ell+1}- F[D_{\ell},k_{\ell}]\right)\nonumber\\
%
&=&\prod_{\ell=0}^{L-1}\sum_{k_{\ell}=0}^T\Prob\left[k_\ell\Bigg\vert D_{0}\prod_{\tilde{\ell}=0}^{\ell-1}\left\{1- h\left[f- 1+\frac{k_{\tilde{\ell}}}{T}\right]\right\}\right]\nonumber\\
&&\times\prod_{\ell=0}^{L-1}\delta\left(D_{\ell+1}- D_{0}\prod_{\tilde{\ell}=0}^{\ell}\left\{1- h\left[f- 1+\frac{k_{\tilde{\ell}}}{T}\right]\right\}\right)
%
\end{eqnarray}
%
Thus
%
\begin{eqnarray}
&&\Prob\left[D_{L},\ldots,D_{1}\vert D_{0}\right]\nonumber\\
&&~~~~~~~~=\prod_{\ell=0}^{L-1}\sum_{k_{\ell}=0}^T\Prob\left[k_\ell\Bigg\vert D_{0}\prod_{\tilde{\ell}=0}^{\ell-1}\left\{1- h\left[f- 1+\frac{k_{\tilde{\ell}}}{T}\right]\right\}\right]\nonumber\\
&&~~~~~~~~~~~~~\times\prod_{\ell=0}^{L-1}\delta\left(D_{\ell+1}- D_{0}\prod_{\tilde{\ell}=0}^{\ell}\left\{1- h\left[f- 1+\frac{k_{\tilde{\ell}}}{T}\right]\right\}\right).
%
\end{eqnarray}
%
From above follows that for average 
%
\begin{eqnarray}
\langle D_{L}\rangle&=&\int\rmd D_{L}\cdots \int\rmd D_{1} \Prob\left[D_{L},\ldots,D_{1}\vert D_{0}\right]\,D_{L}
%
\end{eqnarray}
%
we obtain 
%
\begin{eqnarray}
%
\langle D_{L}\rangle&=&\prod_{\ell=0}^{L-1}\sum_{k_{\ell}=0}^T\Prob\left[k_\ell\Bigg\vert D_{0}\prod_{\tilde{\ell}=0}^{\ell-1}\left\{1- h\left[f- 1+\frac{k_{\tilde{\ell}}}{T}\right]\right\}\right]\nonumber\\
&&\times D_{0}\prod_{\ell=0}^{L-1}\left\{1- h\left[f- 1+\frac{k_{\ell}}{T}\right]\right\}\label{eq:aver-D}.
%
\end{eqnarray}
%
and
%
\begin{eqnarray}
%
\langle D_{L+1}\rangle&=&\prod_{\ell=0}^{L}\sum_{k_{\ell}=0}^T\Prob\left[k_\ell\Bigg\vert D_{0}\prod_{\tilde{\ell}=0}^{\ell-1}\left\{1- h\left[f- 1+\frac{k_{\tilde{\ell}}}{T}\right]\right\}\right]\nonumber\\
&&\times D_{0}\prod_{\ell=0}^{L}\left\{1- h\left[f- 1+\frac{k_{\ell}}{T}\right]\right\}.
%
\end{eqnarray}
%
Consider 
%
\begin{eqnarray}
%
\langle D_{L+1}\rangle&=&\prod_{\ell=0}^{L}\sum_{k_{\ell}=0}^T\Prob\left[k_\ell\Bigg\vert D_{0}\prod_{\tilde{\ell}=0}^{\ell-1}\left\{1- h\left[f- 1+\frac{k_{\tilde{\ell}}}{T}\right]\right\}\right]\nonumber\\
&&\times D_{0}\prod_{\ell=0}^{L}\left\{1- h\left[f- 1+\frac{k_{\ell}}{T}\right]\right\}\nonumber\\
%
&=&\prod_{\ell=0}^{L}\sum_{k_{\ell}=0}^T\Prob\left[k_\ell\Bigg\vert D_{0}\prod_{\tilde{\ell}=0}^{\ell-1}\left\{1- h\left[f- 1+\frac{k_{\tilde{\ell}}}{T}\right]\right\}\right]\nonumber\\
&&\times D_{0}\left\{\prod_{\ell=0}^{L-1}\left\{1- h\left[f- 1+\frac{k_{\ell}}{T}\right]\right\}\right\}\left\{1- h\left[f- 1+\frac{k_{L}}{T}\right]\right\}\nonumber\\
%
&=&\langle D_{L}\rangle- h\langle D_{L}\rangle(f- 1)\nonumber\\
%
&&-h\prod_{\ell=0}^{L}\sum_{k_{\ell}=0}^T\Prob\left[k_\ell\Bigg\vert D_{0}\prod_{\tilde{\ell}=0}^{\ell-1}\left\{1- h\left[f- 1+\frac{k_{\tilde{\ell}}}{T}\right]\right\}\right]\nonumber\\
%
&&\times D_{0}\left\{\prod_{\ell=0}^{L-1}\left\{1- h\left[f- 1+\frac{k_{\ell}}{T}\right]\right\}\right\}\frac{k_L}{T}
%
\end{eqnarray}
%
Now consider the average
%
\begin{eqnarray}
%
&&\prod_{\ell=0}^{L}\sum_{k_{\ell}=0}^T\Prob\left[k_\ell\Bigg\vert D_{0}\prod_{\tilde{\ell}=0}^{\ell-1}\left\{1- h\left[f- 1+\frac{k_{\tilde{\ell}}}{T}\right]\right\}\right]\nonumber\\
%
&&\times D_{0}\left\{\prod_{\ell=0}^{L-1}\left\{1- h\left[f- 1+\frac{k_{\ell}}{T}\right]\right\}\right\}\frac{k_L}{T}\nonumber\\
%
&&=\prod_{\ell=0}^{L-1}\sum_{k_{\ell}=0}^T\Prob\left[k_\ell\Bigg\vert D_{0}\prod_{\tilde{\ell}=0}^{\ell-1}\left\{1- h\left[f- 1+\frac{k_{\tilde{\ell}}}{T}\right]\right\}\right]\nonumber\\
%
&&\times D_{0}\prod_{\ell=0}^{L-1}\left\{1- h\left[f- 1+\frac{k_{\ell}}{T}\right]\right\}\nonumber\\
%
&&\times\left[1-\phi\left(\frac{\sum_{i=1}^Nw_i}{D_{0}\prod_{\ell=0}^{L-1}\left\{1- h\left[f- 1+\frac{k_{\ell}}{T}\right]\right\}}\right)\right]\nonumber\\
%
\end{eqnarray}
%
We note that the function $D(1-\phi(\sum_{i=1}^Nw_i/D))$ is convex in $D$. Hence the average 
%
\begin{eqnarray}
%
&&\prod_{\ell=0}^{L-1}\sum_{k_{\ell}=0}^T\Prob\left[k_\ell\Bigg\vert D_{0}\prod_{\tilde{\ell}=0}^{\ell-1}\left\{1- h\left[f- 1+\frac{k_{\tilde{\ell}}}{T}\right]\right\}\right]\nonumber\\
%
&&~~~~~\times D_{0}\prod_{\ell=0}^{L-1}\left\{1- h\left[f- 1+\frac{k_{\ell}}{T}\right]\right\}\nonumber\\
%
&&~~~~~~~~~\times\left[1-\phi\left(\frac{\sum_{i=1}^Nw_i}{D_{0}\prod_{\ell=0}^{L-1}\left\{1- h\left[f- 1+\frac{k_{\ell}}{T}\right]\right\}}\right)\right]\nonumber\\
%
&&~~~~~~~~~~~~~~~~~~~~~~~~~~~\geq \langle D_{L}\rangle\left[1-\phi\left(\frac{\sum_{i=1}^Nw_i}{\langle D_{L}\rangle}\right)\right]
%
\end{eqnarray}
%


Let us consider the product 
%
%
\begin{eqnarray}
%
\prod_{\ell=0}^{L-1}\left\{1- h\left[f- 1+\frac{k_{\ell}}{T}\right]\right\}&=&\rme^{\sum_{\ell=0}^{L-1}\log\left(1- h\left[f- 1+\frac{k_{\ell}}{T}\right]\right)}
%
\end{eqnarray}

\section{Mean-field approximation and concentration}

Let
\begin{equation}
W=\sum_{i=1}^N w_i,
\qquad
\phi(u)=1-(1-f)^u .
\label{eq-app:phi_def}
\end{equation}

Consider the stochastic recurrence
\begin{equation}
D_{\ell+1}
=
D_\ell-hD_\ell
\left(
f-1+\frac{k_\ell}{T}
\right),
\label{eq-app:stochastic_recursion}
\end{equation}
where, conditionally on $D_\ell$,
\begin{equation}
k_\ell
\sim
\mathrm{Bin}(T,p(D_\ell)),
\label{eq-app:binomial_def}
\end{equation}
with
\begin{equation}
p(D)=1-\phi(W/D).
\label{eq-app:p_def}
\end{equation}

Define the deterministic mean-field recursion
\begin{equation}
d_{\ell+1}
=
d_\ell-hd_\ell
\left(
f-\phi(W/d_\ell)
\right),
\label{eq-app:meanfield_recursion}
\end{equation}
with deterministic initial condition
\begin{equation}
d_0=D_0>0.
\label{eq-app:init_cond}
\end{equation}

Define also
\begin{equation}
g(D)=D-hD(f-\phi(W/D)).
\label{eq-app:g_def}
\end{equation}

Assume
\begin{equation}
h<\frac1f.
\label{eq-app:h_assumption}
\end{equation}

Let $L\ge1$ be fixed.

For the stochastic recurrence (\ref{eq-app:stochastic_recursion}), the extremal updates occur when
\begin{equation}
k_\ell=0
\label{eq:k_zero}
\end{equation}
or
\begin{equation}
k_\ell=T.
\label{eq-app:k_T}
\end{equation}

Since
\begin{equation}
0\le \frac{k_\ell}{T}\le 1,
\label{eq-app:k_bounds}
\end{equation}
equation (\ref{eq-app:stochastic_recursion})  implies
\begin{equation}
1-hf
\le
1-h\left(f-1+\frac{k_\ell}{T}\right)
\le
1+h(1-f).
\label{eq-app:update_factor_bounds}
\end{equation}

Hence
\begin{equation}
D_\ell(1-hf)
\le
D_{\ell+1}
\le
D_\ell(1+h(1-f)).
\label{eq-app:D_step_bounds}
\end{equation}

Iterating (\ref{eq-app:D_step_bounds}) gives
\begin{equation}
D_0(1-hf)^\ell
\le
D_\ell
\le
D_0(1+h(1-f))^\ell.
\label{eq-app:D_global_bounds}
\end{equation}

Since (\ref{eq-app:h_assumption}) implies
\begin{equation}
1-hf>0,
\label{eq-app:positive_factor}
\end{equation}
we may define
\begin{equation}
D_{\min}
:=
D_0(1-hf)^L
>0,
\label{eq-app:Dmin}
\end{equation}
and
\begin{equation}
D_{\max}
:=
D_0(1+h(1-f))^L.
\label{eq-app:Dmax}
\end{equation}

Therefore,
\begin{equation}
D_\ell\in[D_{\min},D_{\max}],
\qquad
0\le \ell\le L.
\label{eq-app:D_interval}
\end{equation}

Similarly, since
\begin{equation}
0\le \phi(W/d_\ell)\le 1,
\label{eq-app:phi_bounds}
\end{equation}
the deterministic recursion satisfies
\begin{equation}
d_\ell(1-hf)
\le
d_{\ell+1}
\le
d_\ell(1+h(1-f)),
\label{eq-app:d_step_bounds}
\end{equation}
hence
\begin{equation}
d_\ell\in[D_{\min},D_{\max}],
\qquad
0\le \ell\le L.
\label{eq-app:d_interval}
\end{equation}

Define the interval
\begin{equation}
I=[D_{\min},D_{\max}].
\label{eq-app:I_def}
\end{equation}

Assume furthermore that $g$ is Lipschitz on $I$ with constant
\begin{equation}
L_g
=
\sup_{D\in I}|g'(D)|
<
1.
\label{eq-app:Lg_assumption}
\end{equation}

The condition $L_g<1$ means that the deterministic mean-field dynamics is locally contractive on $I$. In particular, if
\begin{equation}
x_{\ell+1}=g(x_\ell),
\qquad
y_{\ell+1}=g(y_\ell),
\label{eq-app:two_traj}
\end{equation}
then
\begin{equation}
|x_{\ell+1}-y_{\ell+1}|
\le
L_g|x_\ell-y_\ell|,
\label{eq-app:contraction}
\end{equation}
and therefore
\begin{equation}
|x_\ell-y_\ell|
\le
L_g^\ell |x_0-y_0|.
\label{eq-app:geom_decay}
\end{equation}

Thus perturbations decay geometrically under the deterministic dynamics. This property is essential because stochastic fluctuations are injected at every iteration through the binomial noise. The contractive condition $L_g<1$ ensures that previously generated fluctuations decay rather than accumulate.

\begin{prop}
Assume (\ref{eq-app:h_assumption}) and (\ref{eq-app:Lg_assumption}). Then, for every fixed horizon $L$,
\begin{equation}
\sup_{0\le \ell\le L}|D_\ell-d_\ell|
=
O_{\mathbb P}(T^{-1/2}).
\label{eq-app:main_concentration}
\end{equation}

Furthermore, defining
\begin{equation}
m_\ell:=\mathbb E[D_\ell],
\label{eq-app:m_def}
\end{equation}
we have
\begin{equation}
m_\ell-d_\ell
=
O(T^{-1/2}),
\qquad
0\le \ell\le L.
\label{eq-app:mean_concentration}
\end{equation}
\end{prop}

\begin{proof}

Define the fluctuation
\begin{equation}
\epsilon_\ell=D_\ell-d_\ell.
\label{eq-app:e_def}
\end{equation}

Using
\begin{equation}
\frac{k_\ell}{T}
=
p(D_\ell)+\eta_\ell,
\label{eq-app:eta_decomp}
\end{equation}
where
\begin{equation}
\eta_\ell
=
\frac{k_\ell}{T}-p(D_\ell),
\label{eq-app:eta_def}
\end{equation}
substituting (\ref{eq-app:eta_decomp}) into
\ref{eq-app:stochastic_recursion} gives
\begin{equation}
D_{\ell+1}
=
D_\ell
-
hD_\ell
\left(
f-1+p(D_\ell)+\eta_\ell
\right).
\label{eq-app:substituted}
\end{equation}

Using (\ref{eq-app:p_def}),
\begin{equation}
f-1+p(D)
=
f-\phi(W/D).
\label{eq-app:simplify_p}
\end{equation}

Hence
\begin{equation}
D_{\ell+1}
=
D_\ell
-
hD_\ell
\left(
f-\phi(W/D_\ell)
\right)
-
hD_\ell\eta_\ell.
\label{eq-app:stochastic_eta}
\end{equation}

Using (\ref{eq-app:g_def}),
\begin{equation}
D_{\ell+1}
=
g(D_\ell)-hD_\ell\eta_\ell.
\label{eq-app:g_eta}
\end{equation}

Subtracting the deterministic recursion
(\ref{eq-app:meanfield_recursion}) gives
\begin{equation}
e_{\ell+1}
=
g(D_\ell)-g(d_\ell)
-
hD_\ell\eta_\ell.
\label{eq-app:e_recursion}
\end{equation}

Since $g$ is Lipschitz on $I$,
\begin{equation}
|g(D_\ell)-g(d_\ell)|
\le
L_g|\epsilon_\ell|.
\label{eq-app:g_lipschitz}
\end{equation}

Using also $D_\ell\le D_{\max}$,
equation (\ref{eq-app:e_recursion}) yields
\begin{equation}
|e_{\ell+1}|
\le
L_g|\epsilon_\ell|
+
hD_{\max}|\eta_\ell|.
\label{eq-app:e_ineq}
\end{equation}

Since
\begin{equation}
e_0=D_0-d_0=0,
\label{eq-app:e0}
\end{equation}
iterating (\ref{eq-app:e_ineq}) gives
\begin{equation}
|\epsilon_\ell|
\le
hD_{\max}
\sum_{r=0}^{\ell-1}
L_g^{\ell-1-r}
|\eta_r|.
\label{eq-app:e_sum}
\end{equation}

Because
\begin{equation}
k_\ell\mid D_\ell
\sim
\mathrm{Bin}(T,p(D_\ell)),
\label{eq-app:binomial_repeat}
\end{equation}
we have
\begin{equation}
\mathbb E[\eta_\ell\mid D_\ell]=0,
\label{eq-app:eta_mean}
\end{equation}
and
\begin{equation}
\mathrm{Var}(\eta_\ell\mid D_\ell)
=
\frac{p(D_\ell)(1-p(D_\ell))}{T}.
\label{eq-app:eta_var}
\end{equation}

Since
\begin{equation}
0\le p(D_\ell)\le 1,
\label{eq-app:p_bounds}
\end{equation}
the elementary inequality
\begin{equation}
p(1-p)\le \frac14
\label{eq-app:pineq}
\end{equation}
implies
\begin{equation}
\mathrm{Var}(\eta_\ell\mid D_\ell)
\le
\frac{1}{4T}.
\label{eq-app:eta_var_bound}
\end{equation}

Also,
\begin{equation}
|\eta_\ell|\le 1.
\label{eq-app:eta_abs}
\end{equation}

Define
\begin{equation}
a_{\ell,r}
=
hD_{\max}L_g^{\ell-1-r},
\qquad
r=0,\ldots,\ell-1.
\label{eq-app:a_def}
\end{equation}

Then (\ref{eq-app:e_sum}) becomes
\begin{equation}
|\epsilon_\ell|
\le
\sum_{r=0}^{\ell-1}
a_{\ell,r}|\eta_r|.
\label{eq-app:e_sum2}
\end{equation}

Squaring both sides gives
\begin{equation}
|\epsilon_\ell|^2
\le
\left(
\sum_{r=0}^{\ell-1}
a_{\ell,r}|\eta_r|
\right)^2.
\label{eq-app:square}
\end{equation}

Applying Cauchy--Schwarz,
\begin{equation}
\left(
\sum_{r=0}^{\ell-1}
a_{\ell,r}|\eta_r|
\right)^2
\le
\left(
\sum_{r=0}^{\ell-1}a_{\ell,r}^2
\right)
\left(
\sum_{r=0}^{\ell-1}\eta_r^2
\right).
\label{eq-app:CS}
\end{equation}

Therefore
\begin{equation}
|\epsilon_\ell|^2
\le
\left(
\sum_{r=0}^{\ell-1}a_{\ell,r}^2
\right)
\left(
\sum_{r=0}^{\ell-1}\eta_r^2
\right).
\label{eq-app:e_square_bound}
\end{equation}

Now
\begin{equation}
\sum_{r=0}^{\ell-1}a_{\ell,r}^2
=
h^2D_{\max}^2
\sum_{r=0}^{\ell-1}
L_g^{2(\ell-1-r)}.
\label{eq-app:a_sum}
\end{equation}

Changing variable
\begin{equation}
m=\ell-1-r,
\label{eq-app:m_change}
\end{equation}
gives
\begin{equation}
\sum_{r=0}^{\ell-1}a_{\ell,r}^2
=
h^2D_{\max}^2
\sum_{m=0}^{\ell-1}L_g^{2m}.
\label{eq-app:a_sum2}
\end{equation}

Since $\ell\le L$,
\begin{equation}
\sum_{m=0}^{\ell-1}L_g^{2m}
\le
\sum_{m=0}^{L-1}L_g^{2m}
\le
\sum_{m=0}^{\infty}L_g^{2m}
=
\frac{1}{1-L_g^2}.
\label{eq-app:geom_sum}
\end{equation}

Therefore
\begin{equation}
\sum_{r=0}^{\ell-1}a_{\ell,r}^2
\le
\frac{h^2D_{\max}^2}{1-L_g^2}.
\label{eq-app:a_uniform}
\end{equation}

Furthermore,
\begin{equation}
\mathbb E[\eta_r^2]
=
\mathrm{Var}(\eta_r)
=
O(T^{-1}),
\label{eq-app:eta_second}
\end{equation}
hence
\begin{equation}
\sum_{r=0}^{\ell-1}\mathbb E[\eta_r^2]
\le
\sum_{r=0}^{L-1}\mathbb E[\eta_r^2]
=
O(L/T).
\label{eq-app:eta_sum}
\end{equation}

Since $L$ is fixed,
\begin{equation}
O(L/T)=O(T^{-1}).
\label{eq-app:L_fixed}
\end{equation}

Taking expectation in (\ref{eq-app:e_square_bound}) gives
\begin{equation}
\mathbb E[\epsilon_\ell^2]
=
O(T^{-1}).
\label{eq-app:e_second}
\end{equation}

Applying Chebyshev's inequality yields
\begin{equation}
\epsilon_\ell
=
O_{\mathbb P}(T^{-1/2}).
\label{eq-app:e_prob}
\end{equation}

Since the horizon $L$ is fixed,
\begin{equation}
\sup_{0\le \ell\le L}|\epsilon_\ell|
=
O_{\mathbb P}(T^{-1/2}).
\label{eq-app:e_sup}
\end{equation}

Using \ref{eq-app:e_def},
\begin{equation}
\sup_{0\le \ell\le L}|D_\ell-d_\ell|
=
O_{\mathbb P}(T^{-1/2}),
\label{eq-app:final_main}
\end{equation}
which proves (\ref{eq-app:main_concentration}).

Finally, from \ref{eq-app:e_second},
\begin{equation}
\mathbb E|\epsilon_\ell|
\le
\sqrt{\mathbb E[\epsilon_\ell^2]}
=
O(T^{-1/2}),
\label{eq-app:e_first}
\end{equation}
where the inequality follows from the Cauchy--Schwarz inequality
\[
|\mathbb E[XY]|^2
\le
\mathbb E[X^2]\mathbb E[Y^2]
\]
applied with $X=|\epsilon_\ell|$ and $Y=1$.

Since
\begin{equation}
m_\ell-d_\ell
=
\mathbb E[\epsilon_\ell],
\label{eq-app:m_minus_d}
\end{equation}
we obtain
\begin{equation}
|m_\ell-d_\ell|
\le
\mathbb E|\epsilon_\ell|
=
O(T^{-1/2}),
\label{eq-app:final_mean}
\end{equation}
which proves \ref{eq-app:mean_concentration}.

\end{proof}


%---------------------------------------------------------------------
\section*{References}

%\bibliographystyle{unsrt} 

%\bibliography{ppos}

\end{document}
